\documentclass[11pt,landscape,a4paper]{article}
\input{./preamble}
\input{./color-definition}

%\usepackage[osf]{mathpazo} 	% Palatino with smallcaps and oldstyle numbers, 
\usepackage[sc]{mathpazo}	% altrimenti si usa '[sc]' per il soli smallcaps

% -----------------------------------------------------------------------

\hypersetup{pdfpagemode=FullScreen}

\begin{document}

\raggedright
\footnotesize
\begin{multicols}{3}

% multicol parameters
% These lengths are set only within the two main columns
%\setlength{\columnseprule}{0.25pt}
\setlength{\premulticols}{1pt}
\setlength{\postmulticols}{1pt}
\setlength{\multicolsep}{1pt}
\setlength{\columnsep}{2pt}

%%%%%%%%%%%%%%%%%%%%%%%%%%%%%%%%%%%%%%%%%%%%%%%%%%%%%%%%%%%
%%%%%%%%%%%%%%%%%%%%%%%%%%%%%%%%%%%%%%%%%%%%%%%%%%%%%%%%%%%
%%%%%%%%%%%%%%%%%%%%%%%%%%%%%%%%%%%%%%%%%%%%%%%%%%%%%%%%%%%

\begin{center}
     \Large{\textbf{LM Emacs reference card}} \\
     (2026-04-00 --- keybinding first)
\end{center}

%{\color{darkred} \verb|C-c C-x|} 
%{\color{bblue} \verb|C-c C-f|} 
%{\color{cblue} \verb|C-c C-w|} 
%{\color{cred} \verb|C-c C-w|} 

% semantica dei colori
%
%						verde: richiede package;

%\section{starting Emacs}
%To enter GNU Emacs 27, just type its name: \verb!emacs!

{\itshape This reference card shows the commands and keybindings LM added to his Emacs, they are grouped by the action they made (insert strings, delete objects \dots).}\hfill --- {\itshape\scriptsize \today}

\vspace{1ex}
\hline
\vspace{1.5ex}
%\hline
%\vspace{1em}

\section{visual operations (\textit{windows} / \textit{frame})}

\begin{tabular}{@{}ll@{}}
{\color{darkred}\verb|C-c +|} 			& zoom in\\
{\color{darkred}\verb|C-c -|} 			& zoom out\\
{\color{darkred}\verb|C-c v|} 			& activate/deactivate visual-line-mode\\ & (i.e. word wrap)\\
\\
{\verb|C-l|} 							& redraw garbaged screen\\
{\verb|C-c 1|} or {\color{orange}\verb|F1|}				& delete all other windows\\
{\verb|C-c 2|} 							& split window, above and below\\
{\verb|C-c 3|} 							& split window, side by side\\
{\verb|C-c 0|} 							& delete this window\\
{\verb|C-x o|} or {\color{orange} \verb|F12|}			& cycle windows\\
{\verb|C-M-v|} 							& scroll other window\\
%\verb|M-x compare-windows| & find next place where text in\\
%& selected window does not match\\
%& text in the next window\\
\verb|C-x 4| & prefix for select another window\\ 
\verb|C-x 5| & prefix for select another frame\\ 
{\verb|C-x 5 2|}	& create a new frame (default parameters)\\
{\verb|C-x 5 c|}	& clone frame (same parameters)\\
%C-x 5 b bufname RET
%Select buffer bufname in another frame. This runs switch-to-buffer-other-frame.
%C-x 5 f filename RET
%Visit file filename and select its buffer in another frame. This runs find-file-other-frame. 
{\verb|C-x 5 d| \textit{directory} \verb|RET|}	& \verb|Dired| for directory in another frame\\
%Select a Dired buffer for directory <directory> in another frame. This runs dired-other-frame.
{\verb|C-x 5 r| \textit{filename} \verb|RET|}	& visit file read-only in another frame\\
%Visit file filename read-only, and select its buffer in another frame.



%& (splitting the window if there is only one) \\
%& and select a buffer in that window.\\
%C-x 4 is a prefix key for commands that select another window (splitting the window if there is only one) and select a buffer in that window. Different C-x 4 commands have different ways of finding the buffer to select.
%%%%
%Emacs refers to gui windows as 'frames', and the partitions with a frame are called 'windows'.

%You can open a file in a new frame with C-x 5 f. You can open a file in a different window in the same frame with C-x 4 f.
{\verb|C-x 4 f|}	& open a file in a different window\\
{\verb|C-x 5 f|}	& open a file in a different frame\\
\end{tabular}

%Displaying in Another Window
%
%C-x 4 is a prefix key for commands that select another window (splitting the window if there is only one) and select a buffer in that window. Different C-x 4 commands have different ways of finding the buffer to select.
%
%C-x 4 b bufname RET
%Select buffer bufname in another window. This runs switch-to-buffer-other-window.
%C-x 4 C-o bufname RET
%Display buffer bufname in another window, but don't select that buffer or that window. This runs display-buffer.
%C-x 4 f filename RET
%Visit file filename and select its buffer in another window. This runs find-file-other-window. See section Visiting Files.
%[...]
%C-x 4 .
%Find a tag in the current tags table, in another window. This runs find-tag-other-window, the multiple-window variant of M-. (see section Tags Tables).
%C-x 4 r filename RET
%Visit file filename read-only, and select its buffer in another window. This runs find-file-read-only-other-window. See section Visiting Files.
%%
%%
%C-x 4 f
%Visit a file, in another window (find-file-other-window). Don't alter what is displayed in the selected window.
%C-x 5 f
%Visit a file, in a new frame (find-file-other-frame). Don't alter what is displayed in the selected frame.
%M-x find-file-literally
%Visit a file with no conversion of the contents.
%Displaying in Another Window
%
%C-x 4 is a prefix key for commands that select another window (splitting the window if there is only one) and select a buffer in that window. Different C-x 4 commands have different ways of finding the buffer to select.
%
%C-x 4 b bufname RET
%Select buffer bufname in another window. This runs switch-to-buffer-other-window.
%C-x 4 C-o bufname RET
%Display buffer bufname in another window, but don't select that buffer or that window. This runs display-buffer.
%C-x 4 f filename RET
%Visit file filename and select its buffer in another window. This runs find-file-other-window. See section Visiting Files.
%[...]
%C-x 4 .
%Find a tag in the current tags table, in another window. This runs find-tag-other-window, the multiple-window variant of M-. (see section Tags Tables).
%C-x 4 r filename RET
%Visit file filename read-only, and select its buffer in another window. This runs find-file-read-only-other-window. See section Visiting Files.
%%
%%
%C-x 4 f
%Visit a file, in another window (find-file-other-window). Don't alter what is displayed in the selected window.
%C-x 5 f
%Visit a file, in a new frame (find-file-other-frame). Don't alter what is displayed in the selected frame.
%M-x find-file-literally
%Visit a file with no conversion of the contents.
%Displaying in Another Window
%
%C-x 4 is a prefix key for commands that select another window (splitting the window if there is only one) and select a buffer in that window. Different C-x 4 commands have different ways of finding the buffer to select.
%
%C-x 4 b bufname RET
%Select buffer bufname in another window. This runs switch-to-buffer-other-window.
%C-x 4 C-o bufname RET
%Display buffer bufname in another window, but don't select that buffer or that window. This runs display-buffer.
%C-x 4 f filename RET
%Visit file filename and select its buffer in another window. This runs find-file-other-window. See section Visiting Files.
%[...]
%C-x 4 .
%Find a tag in the current tags table, in another window. This runs find-tag-other-window, the multiple-window variant of M-. (see section Tags Tables).
%C-x 4 r filename RET
%Visit file filename read-only, and select its buffer in another window. This runs find-file-read-only-other-window. See section Visiting Files.
%%
%%
%C-x 4 f
%Visit a file, in another window (find-file-other-window). Don't alter what is displayed in the selected window.
%C-x 5 f
%Visit a file, in a new frame (find-file-other-frame). Don't alter what is displayed in the selected frame.
%M-x find-file-literally
%Visit a file with no conversion of the contents.


\section{visual operations (\textit{buffer})}

\begin{tabular}{@{}ll@{}}
	{\color{darkred} \verb|C-TAB|} 	& switch to next buffer\\
	{\color{darkred} \verb|C-S-TAB|}	& switch to previous buffer\\
	{\color{darkred} \verb|C-x k|} 	& kill buffer, no question\\
	{\color{cblue} \verb|C-x C-b|} 	& 'ibuffer\\ %% in precedenza usavo buffer menu per sostituire il default EMACS
	{\color{darkred}\verb|C-c B|} 	& \verb|'bury-buffer| (\textit{better default})\\
	{\verb|C-x b| \emph{buffer} \verb|RET|} & create a new buffer (in F. mode)\\ %% If no buffer with that name exists, Emacs creates it
	{\verb|M-~|} 		& set a buffer as not modified\\ %% M-x not-modified
	% toggle read-only status of buer C-x C-q - sono uguali??
	%%
	%restore a buffer to its original contents M-x revert-buffer
\end{tabular}

\section{visual operations (\textit{minibuffer / echo area})}

\begin{tabular}{@{}ll@{}}
	{\verb|C-]|} 			& cancel minibuffer if focus is in another window\\ % abort-recursive-edit %% FORSE VA MAPPATO SU F2, in analogia con F1
	{\color{cgreen}\verb|C-f|} 		& Ido: revert to the old-style completion\\
\end{tabular}

\vfill\null

\section{CUA mode}

\begin{tabular}{@{}ll@{}}
	{\color{darkred}\verb|C-c|} 	& copy region												\\ 
	{\color{darkred}\verb|C-v|} 	& paste region											\\ %% conflict with "`scroll next"'
	{\color{darkred}\verb|C-x|} 	& cut region									 			\\
	{\color{darkred}\verb|C-z|} 	& undo last action									\\ %% conflict with suspend!
\end{tabular}

%%%%%%%%%%%%%%%%%%%%%%%%%%%%%%%%%%%%%%%%%%%%%%%%%%
%\vfill\null
\columnbreak
%%%%%%%%%%%%%%%%%%%%%%%%%%%%%%%%%%%%%%%%%%%%%%%%%%


	
\section{editing text operations}


\begin{tabular}{@{}ll@{}}
{\color{darkred}\verb|M-o|} 					&  open newline after point (VI:o)\\ 
{\color{darkred}\verb|M-O|} 					&  open newline before point (VI:O)\\ 
%{\verb|C-o|} 	& open newline after point (\verb|open-line|) \\
%				& separate lines/paragrephs (insert break) (VI:o)\\
%??				& open newline before point (VI:O)\\
\\
\end{tabular}



\begin{tabular}{@{}ll@{}}
	{\color{darkred}\verb|C-c d|} 	& delete/kill whole line/paragraph	\\ 
	{\color{darkred}\verb|C-c y|} 	& copy line/paragraph	(goes in kill ring?)  \\
	{\color{darkred}\verb|C-c C-y|} & duplicate line/paragraph (VI: yyp) \\
	%\\ %% vertical spacing
	{\verb|C-o|} 					& separate lines/paragrephs (insert break) (VI:o)\\
	{\color{darkred}\verb|C-c j|}	& merge/join lines/paragrephs\\
	{\color{darkred}\verb| |} 		& copy backward word\\
\end{tabular}

\begin{tabular}{@{}ll@{}}
	{\color{cblue}\verb|M-SPC|} & cycle-spacing\\ %% EMACS with batteries included
	%% cycles between deleting all whitespace around cursor but one, deleting all whitespace and restoring whitespace
	{\color{cblue}\verb|M-q|} 						&  unfill parargraph\\ 
	{\color{darkred}\verb|M-Q|} 					&  fill parargraph\\ 
	{\color{darkred}\verb|C-c j|} 			& join line 
	%(why \verb|M-^| doesen't works?)\\
\end{tabular}


%Formatting
%
%indent current line (mode-dependent) TAB
%indent region (mode-dependent) C-M-\
%indent sexp (mode-dependent) C-M-q
%indent region rigidly arg columns C-x TAB
%indent for comment M-;
%insert newline after point C-o
%move rest of line vertically down C-M-o
%delete blank lines around point C-x C-o
%join line with previous (with arg, next) M-^
%delete all white space around point M-\
%put exactly one space at point M-SPC
%fill paragraph M-q
%set ll column to arg C-x f
%set prex each line starts with C-x .


\subsection{cycling actions}

\begin{tabular}{@{}ll@{}}
	{\color{cblue}\verb|C-w|} 		& kill (move to kill ring)\\
	{\color{cblue}\verb|C-y|} 		& yank (copy from kill ring)\\
	{\verb|C-z|} 		& undo action (or \verb|C-/| or \verb|C-u| or \verb|C-_|)\\
	{\color{cblue}\verb|M-y|} 		& yank former item in kill ring\\
\end{tabular}

\subsection{case operations}

\begin{tabular}{@{}ll@{}}
	{\color{cblue}\verb|C-x C-l|} 		& lowercase region  (\textit{re-enabled keybindings})\\
	{\color{cblue}\verb|C-x C-u|} 		& uppercase region (\textit{re-enabled keybindings})\\
\end{tabular}

\section{comment operations}
\begin{tabular}{@{}ll@{}}
	{\verb|M-;|} 						& comment region \\ 
	{\verb|C-x C-;|} 						& comment or uncomment line/region\\ 
	%Comment or uncomment the current line (comment-line). If the region is active, comment or uncomment the lines in the region instead.
\end{tabular}


\vfill\null

\section{mark operations (select area)}

\begin{tabular}{@{}ll@{}}
	{\verb|M-h|} 		&  mark paragraph\\ 
	{\color{gray}\verb|C-x h|} & mark entire buffer [???]\\ %% disabled! USE: 
	{\verb|C-SPC|}	&  set mark at where point is\\  %%or \verb|C-@|}
	%% Set mark at where point is, clear mark, or jump to mark.
	%{\color{darkred}\verb|C-x C-x|} 		&  (\verb|exchange-point-and-mark|)\\ 
	%% confligge con C-x!!! %% C-c C-c è libero... 	%% forse bisogna riassegnarlo
	%% il comando exchange-point-and-mark. Scambia la posizione del cursore (point) con quella del "mark" (l'inizio della selezione). È utile per controllare l'inizio della regione evidenziata, estenderla o passare rapidamente da un'estremità all'altra di una selezione.
	{\color{darkred}\verb|C-c x|} & exchange-point-and-mark \\
	& (LM: perché \verb|C-x C-x| non va col CUA mnode attivo)\\
	{\color{orange}\verb|C-=|} 		&  expand region\\ 
\end{tabular}


%%%%%%%%%%%%%%%%%%%%%%%%%%%%%%%%%%%%%%%%%%%%%%%%%%
%\vfill\null
\columnbreak
%%%%%%%%%%%%%%%%%%%%%%%%%%%%%%%%%%%%%%%%%%%%%%%%%%


\section{generate text at point}

\begin{tabular}{@{}ll@{}}
	{\color{darkred}\verb|C-c i r s|} 			&  insert random string\\ 
	{\color{darkred}\verb|C-c i r n|} 			&  insert random number\\ 
	{\color{darkred}\verb|C-c i r h|} 			&  insert random hex number\\ 
	{\color{darkred}\verb|C-c i u|}         &  insert UNIX time\\
	{\color{darkred}\verb|C-c i d|}         &  insert date\\ 
	{\color{darkred}\verb|C-c i o|}         &  insert ordinal date\\ 
	{\color{darkred}\verb|C-c C-f|} 	    &  insert filename (at point)\\ 
	{\color{darkred}\verb|C-c C-t|} 		&  insert LM own timestamp (template)\\ 
	{\color{darkred}\verb|C-c C-n|} 		&  insert LM glossa\\ 
	\\
	{\color{cgreen}\verb|C-c l p|}		&  insert \emph{lorem ipsum} paragraph\\ 
	{\color{cgreen}\verb|C-c l s|}		&  insert \emph{lorem ipsum} sentence\\ 
	{\color{cgreen}\verb|C-c l l|}		&  insert \emph{lorem ipsum} list item\\ 
\end{tabular}

\section{operations on files (LM original)}

\begin{tabular}{@{}ll@{}}
	{\color{darkred}\verb|C-c e|} 			& load file \verb|emacs-book.org| in buffer\\
	{\color{darkred}\verb|C-c n|} 			& load file \verb|notes.txt| in buffer\\
	{\color{darkred}\verb|C-c c|} 			& load file \verb|org-capture.org| in buffer\\
	{\color{darkred}\verb|C-c f|} 			& load file \verb|F-T-G.txt| in buffer\\
	{\color{darkred}\verb|C-c z|} 			& load file \verb|zibalduccio.org| in buffer\\
	{\color{orange}\verb|C-c z|} 			& load file \verb|.org| in buffer\\
\end{tabular}


\vspace{1em}

\section{quick help}

\begin{tabular}{@{}ll@{}}
	{\verb|C-h c| \textit{keybinding}} & help on keybinding (\textit{echo area})\\
	%% displays in the echo area the name of the command that key is bound to
	{\verb|C-h k| \textit{keybinding}} & more help on keybinding (\textit{minibuffer})\\
	%% displays a help buffer containing the command?s documentation string, which describes exactly what the command does.
	{\verb|C-h w| \textit{command} \verb|RET|} &  keys that are bound to command (\textit{echo area})\\
	%%  lists the keys that are bound to command. It displays the list in the echo area. If it says the command is not on any key, that means you must use M-x to run it. C-h w runs the command where-is.
	{\verb|C-h m|} &  list all keybindings in currrent major mode\\
	{\verb|M-x which-key-mode|} & displays keybindings  (\textit{global minor mode})\\
	%% M-x which-key is a global minor mode which helps in discovering keymaps. 
	%% It displays keybindings following your currently entered incomplete command (prefix), in a popup.
	{\verb|C-h v| \textit{varname} \verb|RET|} &  status of variable varname (\textit{echo area})\\
	
\end{tabular}

\section{retrieve operation}

\begin{tabular}{@{}ll@{}}
	{\color{cgreen}\verb|M-n|} 		& move to next symbols under point\\
	{\color{cgreen}\verb|M-p|} 		& move to previous symbols under point\\
	{\color{darkred}\verb|C-c o|} 			& occur in buffer\\
	{\color{cblue}\verb|M-z |\emph{char}} 	& zap to char (excluded)  (\verb|'zap-up-to-char|)\\
	{\color{cblue}\verb|M-Z |\emph{char}}	& zap to char (included)  (\verb|'zap-to-char|)\\
	{\color{cblue}\verb|C-s|} 		& 'isearch-forward-regexp\\  %% from Better default
	{\color{cblue}\verb|C-S-s|} 	& 'isearch-forward\\         %% from Better default
	{\color{cblue}\verb|C-r|} 		& 'isearch-backward-regexp\\ %% from Better default
	{\color{cblue}\verb|C-S-r|} 	& 'isearch-backward\\        %% from Better default
	\\
	{\verb|C-M-s|} 	& regexp search foreward\\
	{\verb|C-M-s|} 	& regexp search backward\\
	{\color{cgreen}\verb|C-c t|} 			& search in offline thesaurus\\ 
	\verb|C-s C-s|& repeat the last search by typing \\
\end{tabular}

\section{replace operation}

\begin{tabular}{@{}ll@{}}
	\verb|M-%| \textit{string} \verb|RET| \textit{newstring} \verb|RET| 		& move to next symbols under point [??]\\
\end{tabular}
% Searches in Emacs are case-insensitive, so long as the search string is composed of lower-case letters only. This applies to normal search functions, search-and-replace functions, and regular expression searches. To make a search case-sensitive, all you have to do is specify a capital letter in the search string.

% Emacs also has several forms of advanced searching, such as word search and regular expression search.

% Word search looks for matches to the specified string while ignoring such non-word interlopers like punctuation, whitespace, and formatting characters. It is useful when you’re editing formatted documents such as LaTeX files, or for when you need to find a string but you don’t know exactly what characters may be in or around it. Starting a word search is much like starting a non-incremental search; type Ctrl-s, then press Enter. Then comes the difference — you type Ctrl-w, then the search string, then finally Enter. So the word search Ctrl-s Enter Ctrl-w DVD RW Enter would find both DVD-RW and DVD+RW.

% The most powerful searching tool in Emacs is a regular expression search. You can match single- and multi-character wildcards, ranges of characters, empty strings, and newlines, all according to the standard regular expression syntax. To start an (incremental) regular expression search, just press Ctrl-Meta-s or Ctrl-Meta-r instead of Ctrl-s and Ctrl-r.

%Regular expression quick reference
%. (period) – matches any single character except a newline.
%* (asterisk) – matches zero or more instances of the preceding character.
%+ (plus) – matches one or more instances of the preceding character.
%? (question mark) – matches zero or one (but not more) instances of the preceding character.
%{n} (backslash-braces) – matches exactly n repetitions of the preceding character.
%[ ... ] (brackets) - delimit a set of possible characters.
%^ (caret) – matches the beginning-of-line character.
%$ (period) – matches the end-of-line character.

%%%%%%%%%%%%%%%%%%%%%%%%%%%%%%%%%%%%%%%%%%%%%%%%%%%%%%%%%%%

%Interactive Find Replace, by Regex Pattern
%Alt+x query-replace-regexp
%
%Find replace by Emacs: Regular Expression
%Works on text selection, if none, start from cursor position to end of buffer.
%
%Find Replace by Regex, with Word Boundary
%Ctrl+u Alt+x query-replace-regexp
%
%this saves you time of typing the regex word boundary chars. e.g. \bsomething\b
%
%Find Replace All, by Regex, in One Shot
%Alt+x replace-regexp
%
%Find replace by Emacs: Regular Expression
%Works on text selection, if none, start from cursor position to end of buffer.

\vfill\null




%%%%%%%%%%%%%%%%%%%%%%%%%%%%%%%%%%%%%%%%%%%%%%%%%%
%\vfill\null
\columnbreak
%%%%%%%%%%%%%%%%%%%%%%%%%%%%%%%%%%%%%%%%%%%%%%%%%%



\section{[F]unction keys}

%% NUOVA IPOTESI DI ORGANIZZAZIONE

\begin{tabular}{@{}ll@{}}
\verb!F1!	  & {\color{orange} delete-other-windows} (\verb|C-x 1| alias)\\
\verb!F2!		& {\color{orange} abort-recursive-edit} (\verb|C-]| alias)\\
\verb!F3!   & start defining a keyboard macro\\ 
\verb!F4!   & end / execute last-defined keyboard macro\\
\\  
\verb!F5!   & {\color{darkred} reload \verb|init.el|}\\ 
\verb!F6!   & {\color{darkred} create empty buffer \verb|*fourtytwo*|}\\ 
\verb!F7!		& {\color{orange}\verb|(kill-other-buffers)|}\\ 
\verb!F8!   & {\color{orange} repeat last command} (\verb|C-x z| alias)\\ 
\\
\verb!F9!		& {\color{orange} \verb|(recentf-open-files)|}\\ 
\verb!F10!	& menu (\emph{no hope!})\\ 
\verb!F11!	& maximixe frame on/off\\ 
\verb!F12!  & {\color{orange} next window} (\verb|C-x o| alias)\\ %% formerly help
\end{tabular}

\vspace{1em}

\section{extended keyboard or SO dependent}

\begin{tabular}{@{}ll@{}}
	\verb!C-INS! & copy region\\ 
	\verb!S-INS! & paste region\\ 
	\verb!S-DEL! & cut region\\ 
		\\
		\verb!C-BCKSP! & delete word after point\\ 
		\verb!C-DEL! & delete word after point\\ 
		\verb!C-HOME! & go to begin of buffer\\ 
		\verb!C-END! & go to end of buffer \\ 
		{\verb|C-UP|} 	& go to backward-paragraph\\
		{\verb|C-DOWN|} & go to forward-paragraph\\ 
		\\
		\verb!HOME! & go to begin of paragraph \\
		\verb!END! & go to end of paragraph \\
		\verb!PGUP! & scrolls up one screen\\
		\verb!PGDN! & scrolls down one screen\\
%Home (<home>): move-beginning-of-line (same as C-a).
%End (<end>): move-end-of-line (same as C-e).
%Page Up (<prior>): scroll-down-command (scrolls up one screen).
%Page Down (<next>): scroll-up-command (scrolls down one screen). 
\end{tabular}

\section{OS operations}

\begin{tabular}{@{}ll@{}}
	\verb!C-A-e! & run EMACS (Windows shortcut)\\ 
%	\verb!S-INS! & paste region\\ 
%	\verb!S-DEL! & cut region\\ 
\end{tabular}

%{\color{darkred} print \verb|"  HelLo woRld!!  "|}
% \verb!F2!		& {\color{darkred} visit \verb|~/note.org|} (\verb|C-c n| alias)\\ 

\vspace{1em}



%%%%%%%%%%%%%%%%%%%%%%%%%%%%%%%%%%%%%%%%%%%%%%%%%%
\vfill\null
\columnbreak
%%%%%%%%%%%%%%%%%%%%%%%%%%%%%%%%%%%%%%%%%%%%%%%%%%

\section{LISP operations}

\begin{tabular}{@{}ll@{}}
	{\color{darkred}\verb|C-c C-j|} & \verb|evaluate-buffer|\\ 
	%% M-x eval-buffer %% ma esiste già!?
\end{tabular}

\subsection{evaluate LISP code}

\begin{tabular}{@{}ll@{}}
	{\verb|C-x C-e|} 					& evaluate expression before point\\
	{\verb|C-j|} 							& in \verb|*scratch*|	buffer		\\
	{\verb|C-M-x|} 					  & in any buffer		\\
	{\verb|C-z|} 							& in \verb|emacs-lisp-mode|\\
	{\verb|M-: (ex-pression)|} 							& eval expression	in minibuffer\\
	{\verb|M-x eval-defun|}		& eval defun containing or after\\
	{\verb|M-x eval-region|}	& eval region\\
	{\verb|M-x eval-buffer|} 	& eval region\\
\end{tabular}

\vspace*{1em}

\section{operations enabled by packages}

\subsection{'lorem-ipsum}


\begin{tabular}{@{}ll@{}}
	{\color{cgreen}\verb|C-c l p|} 		& lorem-ipsum-insert-paragraphs\\
	{\color{cgreen}\verb|C-c l s|} 		& lorem-ipsum-insert-sentences\\
	{\color{cgreen}\verb|C-c l l|} 		& lorem-ipsum-insert-list\\
\end{tabular}

\subsection{'smartscan} % cfr. Mastering Emacs!


\begin{tabular}{@{}ll@{}}
	{\color{cgreen}\verb|M-n|} 		& move next symbols under point\\
	{\color{cgreen}\verb|M-p|} 		& move previous symbols under point\\
	{\color{cgreen}\verb|M-'|} 		& replace all symbols matching the one under point\\ %% all symbols - in buffer -
	{\color{cgreen}\verb|C-u M-'|} 		& replace symbols in your current defun only\\
	
\end{tabular}

%Usage
%
%M-n and M-p move between symbols and type M-' to replace all symbols in the buffer matching the one under point, and C-u M-' to replace symbols in your current defun only (as used by narrow-to-defun.)
%
%For more information on how to use Smart Scan and how to master movement in Emacs, read my article on Effective Editing I: Movement.

%% cfr. https://www.masteringemacs.org/article/effective-editing-movement

\subsection{'columnize}


\begin{tabular}{@{}ll@{}}
	{\color{cgreen}\verb|C-x c|} 		& formats a list of items into columns\\
\end{tabular}

\subsection{'mthesaur}

\begin{tabular}{@{}ll@{}}
	{\color{cgreen}\verb|C-c t|} 			& search in offline thesaurus (\verb|mthesaur-search|)\\ 
\end{tabular}


\subsection{'word-count-mode}

\begin{tabular}{@{}ll@{}}
	{\color{cgreen}\verb|M-+|}  & \verb|word-count-mode|\\
\end{tabular}

\subsection{'re-builder}

\begin{tabular}{@{}ll@{}}
	{\color{cgreen}\verb|M-x re-builder|} & [trovare assolutamente un keybinding!]\\
\end{tabular}



%%%%%%%%%%%%%%%%%%%%%%%%%%%%%%%%%%%%%%%%%%%%%%%%%%
\vfill\null
\columnbreak
%%%%%%%%%%%%%%%%%%%%%%%%%%%%%%%%%%%%%%%%%%%%%%%%%%

\section{hacks (\textit{anche duplicati altrove})}

\begin{tabular}{@{}ll@{}}
	{\verb|M-~|} 		& set a buffer as not modified\\ %% M-x not-modified
	{\verb|C-]|} 			& cancel minibuffer if your focus is in another window\\ % abort-recursive-edit %% FORSE VA MAPPATO SU F2, in analogia con F1
	{\color{cgreen}\verb|M-n|} 		& move next symbols under point\\
	{\color{cgreen}\verb|M-p|} 		& move previous symbols under point\\
	{\color{cblue}\verb|M-SPC|} & cycle-spacing 
	%% EMACS with batteries included
	%% cycles between deleting all whitespace around cursor but one, deleting all whitespace and restoring whitespace
	(cycles between: \\ & deleting all whitespace around cursor \\ & but one, deleting all whitespace\\ & and restoring whitespace)\\
	{\verb|C-x =|} 						& describe the character after point and its position\\ 
	{\verb|C-x RET f|}				& sets the file coding system for the current buffer\\
	{\verb|M-;|} 						& comment region (alternativa a TAB nel modo Elisp)\\ 
	{\verb|C-x C-;|} 						& comment or uncomment line/region\\ 
	\\
\end{tabular}
% The command M-x comment-region is equivalent to calling M-; on an active region,

\begin{tabular}{@{}ll@{}}
\verb|M-x compare-windows|  & find next place where text\\ 
& in selected window does not match\\
&text in next window\\
{\verb|C-x b| \emph{buffer name} \verb|RET|} & create a new buffer (in F. mode)\\ %% If no buffer with that name exists, Emacs creates it
\verb|M-x normal-mode| & revert buffer to F. mode\\
\end{tabular}
%\vspace{2em}


%\section{odds couples} %% oppure missing couples - le comandi che avrebbe senso avere come coppia di keybindings simmetrici

%\section{swaps} %% dove il keybindings originario è stato rimappato con uno nuovo e viceversa (o quasi)

%\section{equivalences/alternatives}

%\section{assegna nuovi chords}


\section{meta operations}

\begin{tabular}{@{}ll@{}}
	\verb|C-x z| or \verb|F8|  		& repeat last command\\ 
	\verb|C-x ESC ESC|  & repeat complex command\\ 
	& that used the minibuffer to get arguments\\
	\verb|M-p| and \verb|M-n| & scan through all the different complex commands\\
	%M-p and M-n (and also up-arrow and down-arrow, if your keyboard has these keys) to scan through all the different complex commands
	
\end{tabular}

%Use the repeat command (C-x z) to repeat the last command. If you preface it with a prefix argument, the prefix arg is applied to the command.
%
%You can also type C-x ESC ESC (repeat-complex-command) to reinvoke commands that used the minibuffer to get arguments. In repeat-complex-command you can type M-p and M-n (and also up-arrow and down-arrow, if your keyboard has these keys) to scan through all the different complex commands you’ve typed.
%
%To repeat a set of commands, use keyboard macros. Use C-x ( and C-x ) to make a keyboard macro that invokes the command and then type C-x e. See Keyboard Macros in The GNU Emacs Manual.


\section{rectangle operations}

% The rectangle operations fall into two classes: commands to erase or insert rectangles, and commands to make blank rectangles.

\begin{tabular}{@{}ll@{}}
	\verb|C-x SPC| 		& toggle rectangle mark mode\\ 
	
\end{tabular}



%C-x r k
%Kill the text of the region-rectangle, saving its contents as the last killed rectangle (kill-rectangle).
%C-x r y
%Yank the last killed rectangle with its upper left corner at point (yank-rectangle).
%
%C-x r o
%Insert blank space to fill the space of the region-rectangle (open-rectangle). This pushes the previous contents of the region-rectangle to the right.
%
%C-x SPC
%Toggle Rectangle Mark mode (rectangle-mark-mode). When this mode is active, the region-rectangle is highlighted and can be shrunk/grown, and the standard kill and yank commands operate on it.
%
%C-x r N
%Insert line numbers along the left edge of the region-rectangle (rectangle-number-lines). This pushes the previous contents of the region-rectangle to the right.
%
%C-x r c
%Clear the region-rectangle by replacing all of its contents with spaces (clear-rectangle).
%
%M-x delete-whitespace-rectangle
%Delete whitespace in each of the lines on the specified rectangle, starting from the left edge column of the rectangle.
%
%C-x r t string RET
%Replace rectangle contents with string on each line (string-rectangle).
%
%M-x string-insert-rectangle RET string RET
%Insert string on each line of the rectangle.

\vfill\null

\section{smista}

\begin{tabular}{@{}ll@{}}
	{\color{darkred}\verb|C-c b|} 	& ding a bell (\verb|ding|) - test keybinding\\
	{\color{darkred}\verb|C-c q|} 	& random quote from \verb|~/.random-lines.txt|\\
	%{\color{darkred}\verb|???|} \verb|F8|		& open an empty buffer (*fourtytwo*)\\
	{\color{darkred}\verb|M-/|} 	& \verb|'hippie-expand| (\textit{better default})\\
	{\verb|M-~|} 		& set a buffer as not modified\\ %% M-x not-modified
	{\verb|M-+|} 		&  word-count-mode\\
	\verb|M-x normal-mode| & revert buffer to F. mode\\
\end{tabular}

\section{smista}

\begin{tabular}{@{}ll@{}}
	{\color{darkred}\verb|C-c C-w g|} 			& search in Google \\
	{\color{darkred}\verb|C-c C-w y|} 			& search in YouTube \\ 
\end{tabular}


%%%%%%%%%%%%%%%%%%%%%%%%%%%%%%%%%%%%%%%%%%%%%%%%%%
%\strut
%\vfill\null
%%%%%%%%%%%%%%%%%%%%%%%%%%%%%%%%%%%%%%%%%%%%%%%%%%
\vspace*{1em}
\hline

\begin{center}
	{\color{bblue}\emph{point - mark - region - rectangle - register - ring(s) - sexp - defun - symbol}}
\end{center}

%% cfr.: https://www.emacswiki.org/emacs/Glossary

% A symbol in GNU Emacs Lisp is an object with a name. The symbol name serves as the printed representation of the symbol.

% The word ‘sexp’ is derived from “s-expression”, the ancient term for an expression in Lisp. But in Emacs, the notion of ‘sexp’ is not limited to Lisp.  

% page – A block of text delimited by form feed (aka page break) characters or the beginning or end of the buffer

% mark – a buffer location that you can change, return to…

% point – where text insertion occurs (sometimes called “cursor position” outside Emacs).

% rectangle – characters between a pair of columns on the screen.

% register – a location for temporarily saving text (a string), a rectangle, a buffer position, a frame or window configuration

%ring – a data structure that acts as if it is circular (no end) 

%% rings: mark, search, kill (local / global)

% kill – cut, that is, delete and copy to the kill ring


\begin{center}
	scope: {\color{bblue}\emph{point - mark | char - word - sentence - line/paragraph - page - whole text in buffer | region}}
\end{center}

\begin{center}
	scope: {\color{bblue}\emph{buffer - windows - minibuffer/echo area - frame - EMACS}}
\end{center}

% frame – Emacs windows are shown in frames (called “windows” outside of Emacs)

% chord – a key sequence with keys pressed simultaneously
% isearch – incremental search, which interactively searches for the search string
% mode line – at the bottom of a window, it describes the current buffer status
% regexp – a regular expression

% The echo area is not a buffer – you cannot edit anything there. 
% It is simply a space for displaying messages. Most text written to the echo area is also logged to the MessagesBuffer.

% To cancel whatever you’re doing in the minibuffer, use ‘C-g’. 
% To cancel minibuffer if your focus is in another window, try the command ‘abort-recursive-edit’ which by default is bound to `C-]’.

%%%%%%%%%%%%%%%%%%%%%%%%%%%%%%%%%%%%%%%%%%%%%%%%%%
%\vfill\null
\columnbreak
%%%%%%%%%%%%%%%%%%%%%%%%%%%%%%%%%%%%%%%%%%%%%%%%%%

\section{reminder}


\begin{tabular}{@{}ll@{}}
	{\verb|M-g M-g|}	 				& \textbf{go to specific line}\\ 
	{\verb|C-l|}	 						& \textbf{scroll point (top centre bottom)}\\ 
	{\verb|M-r|}	 						& \textbf{move point (top centre bottom)}\\ 
	{\verb|C-x z|} 						& \textbf{repeat last command}\\ 
	{\verb|C-x M-z|} 					& (repeat-complex-command)\\
	{\verb|C-M-v|}            & scroll next window upward \\ %(scroll-other-window) %% NONOSTANTE IL CUA mode
	{\verb|C-M-S-v|}          & scroll the next window downward\\ %(scroll-other-window-down) %% NONOSTANTE IL CUA mode
\end{tabular}

%Repositioning Point
%
%You can move the point between the top, center and bottom (by default) of the visible window but without actually scrolling up or down. The command is sometimes useful if you want to reach text in one of the previous three areas, although you can configure the variable recenter-positions to change that. The command is bound to M-r which makes it very easy to reach and thus use.

%Goto Line
%
%Emacs does, of course, support jumping to a specific line, and it is bound to two commands: M-g g (that everybody seems to use), and the easier-to-type M-g M-g.

%%% cfr. https://www.masteringemacs.org/article/effective-editing-movement

\vspace{1em}

%\emph{Qui vanno i comandi standard che non serve abilitare, ma che è opportuno ricordare}

% You can change the encoding to use for the file when saving using ‘C-x C-m f’. You can also force this immediately by using ‘C-x C-m c <encoding> RET C-x C-w RET’. You can list all available encodings with ‘M-x list-coding-systems’.

\begin{tabular}{@{}ll@{}}
	{\verb|C-x C-m f|}		& change encoding for the file\\
	{\verb|C-y|} 					& appends current kill to search string \\ % (\verb|isearch-yank-kill|)
	{\verb|M-y|} 					& replaces that appended text with an earlier kill \\ % (\verb|isearch-yank-pop|) %% if called after C-y, replaces that appended text with an earlier kill
	{\verb|C-w|} 					& appends next char or word at point to search string \\
	{\verb|C-x C-j|} 			& jump to the directory of current buffer (\verb!'dired-x!)\\ 
	{\verb|M-BKSP|} 			& delete the previous word\\
	{\verb|C-s C-s|} 			& repeat last search forward\\
	{\verb|C-r C-r|} 			& repeat last search backward\\
	{\verb|C-x C-x|} 			& exchange point and mark\\
	{\verb|M-SPC|}        & leave single space between words\\
	{\verb|C-x f| \emph{arg}}       & set ll column to arg\\
	%{\verb|C-z|} 							& 									\\ 
\end{tabular}

\vspace{1em}

\begin{tabular}{@{}ll@{}}
	{\verb|C-h h|}	 					& display the HELLO file\\ 
	%{\verb|C-z|} 							& 									\\ 
\end{tabular}

\vspace{1em}

\begin{tabular}{@{}ll@{}}
	{\verb|C-x 3 M-x follow-mode|}	 					& v. split windows and follow mode\\ 
	%{\verb|C-x|} 							& 									\\
	%{\verb|C-z|} 							& 									\\ 
\end{tabular}

\vspace*{1em}
\hline
\vspace*{1em}

Building regexps and replacing constructed regexps: these two commands do complementary things and deserve to work together. As things stand now, you?d have to:

\vspace*{1em}

Open re-builder and construct your regexp.\\
Copy it to the kill ring (C-c C-w by default).\\
Quit re-builder (C-c C-q by default).\\
Run query-replace-regexp (C-M-\% by default).\\
Paste the kill (C-y).\\
Fix newlines and backslashes (if using the ?read? interface to re-builder, more on this below)\\
Type in the replacement string and press RET to begin the replacements.

\vspace*{1em}

Connecting commands through the kill-ring (or clipboard) should be the last resort, not a go-to strategy!


%%%%%%%%%%%%%%%%%%%%%%%%%%%%%%%%%%%%%%%%%%%%%%%%%%
\vfill\null
\columnbreak
%%%%%%%%%%%%%%%%%%%%%%%%%%%%%%%%%%%%%%%%%%%%%%%%%%

\section{LM keybindings check}

\begin{tabular}{@{}ll@{}}
	{\color{darkred}\verb|C-c a|} & \\
	{\color{darkred}\verb|C-c b|} & \\
	{\color{darkred}\verb|C-c c|} & \\
	{\color{darkred}\verb|C-c d|} & \\
	{\color{darkred}\verb|C-c e|} & \\
	{\color{darkred}\verb|C-c f|} & \\
	{\color{darkred}\verb|C-c g|} & \\
	{\color{darkred}\verb|C-c h|} & \\
	{\color{darkred}\verb|C-c i|} & \\
	{\color{darkred}\verb|C-c l|} & \\
	{\color{darkred}\verb|C-c m|} & \\
	{\color{darkred}\verb|C-c n|} & \\
	{\color{darkred}\verb|C-c o|} & \\
	{\color{darkred}\verb|C-c p|} & \\
	{\color{darkred}\verb|C-c q|} & \\
	{\color{darkred}\verb|C-c r|} & \\
	{\color{darkred}\verb|C-c s|} & \\
	{\color{darkred}\verb|C-c t|} & \\
	{\color{darkred}\verb|C-c u|} & \\
	{\color{darkred}\verb|C-c v|} & \\
	{\color{darkred}\verb|C-c z|} & \\
	{\color{darkred}\verb|C-c j|} & \\
	{\color{darkred}\verb|C-c k|} & \\
	{\color{darkred}\verb|C-c w|} & \\
	{\color{darkred}\verb|C-c x|} & \\
	{\color{darkred}\verb|C-c y|} & \\
	\\
	{\color{darkred}\verb|C-c ~|} & \\
	{\color{darkred}\verb|C-c !|} & \\
	{\color{darkred}\verb|C-c @|} & \\
	{\color{darkred}\verb|C-c #|} & \\
	{\color{darkred}\verb|C-c $|} & \\
	{\color{darkred}\verb|C-c %|} & \\
	{\color{darkred}\verb|C-c ^|} & \\
	{\color{darkred}\verb|C-c &|} & \\
	{\color{darkred}\verb|C-c *|} & \\
	{\color{darkred}\verb|C-c (|} & \\
	{\color{darkred}\verb|C-c )|} & \\
	{\color{darkred}\verb|C-c _|} & \\
	{\color{darkred}\verb|C-c +|} & \\
\end{tabular}

%Cerca i modelli. Per esempio, tutto ciò che è sotto C-x C-k riguarda le macro da tastiera. Tutto ciò che è sotto C-x a riguarda le abbreviazioni. E tutto ciò che è sotto C-x r riguarda rettangoli, registri o segnalibri (che sono simili come concetto ai registri).
%
%Inoltre, M-x which-key-mode.

%Unconditional Replacement
%
%M-x replace-string RET string RET newstring RET
%Replace every occurrence of string with newstring.
%M-x replace-regexp RET regexp RET newstring RET
%Replace every match for regexp with newstring.
%To replace every instance of `foo' after point with `bar', use the command M-x replace-string with the two arguments `foo' and `bar'. Replacement happens only in the text after point, so if you want to cover the whole buffer you must go to the beginning first. All occurrences up to the end of the buffer are replaced; to limit replacement to part of the buffer, narrow to that part of the buffer before doing the replacement (see section Narrowing).
%
%When replace-string exits, point is left at the last occurrence replaced. The position of point where the replace-string command was issued is remembered on the mark ring; use C-u C-SPC to move back there.
%
%A numeric argument restricts replacement to matches that are surrounded by word boundaries.

\vfill\null

\section{no keybindings at all (pkos.el)} 

% attivazione a tempo
\begin{tabular}{@{}ll@{}}
\verb|imenu|  &  structural index jump navigation system\\
% imenu — Built-in structural index and jump navigation system for quickly moving through functions, headings, and major sections in large buffers.
	
\verb|which-key|  & available keybindings as you type\\
% displays available keybindings in popup form as you type partial key sequences
\verb|repeat-mode|  & allows rapid repetition of related commands\\
%(repeat-mode 1) ; allows rapid repetition of related commands ; using shorter key sequences, after the initial command
\verb|global-auto-revert|  & refresh buffers when files change on disk\\
\end{tabular}

%%%%%%%%%%%%%%%%%%%%%%%%%%%%%%%%%%%%%%%%%%%%%%%%%%
%\vfill\null
\columnbreak
%%%%%%%%%%%%%%%%%%%%%%%%%%%%%%%%%%%%%%%%%%%%%%%%%%

\section{no keybindings functions}

\begin{tabular}{@{}ll@{}}
{\verb|M-x replace-regexp|} & \\
\verb|M-x compare-windows|  & find next place where text \\ & in selected window does not\\ &  match text in next window\\
{\verb|M-x not-modified|} & set a buffer as not modified (\verb|M-~|)\\
{\verb|M-x join-line|} & join line (\verb|M-^|)\\ %% aka "delete-indentation" % LM keybinding C-c j
{\verb|M-x shell|} & \\
{\verb|M-x scroll-all-mode|} & \\ %% sincronizza due buffer *diversi* NB: va in conflitto col CUA mode!!
{\verb|M-x flyspell-mode|} & \\
{\color{cgreen}\verb|M-x wc-mode|} & word count minor mode\\
{\verb|M-x revert-buffer|} & \\ %% NO KEYBINDINGS: too dangerous!
{\verb|M-x recover-session|} & \\
{\verb|M-x append-to-buffer|} & \\
{\verb|M-x prepend-to-buffer|} & \\
{\verb|M-x append-to-file|} & \\
{\verb|M-x prepend-to-file|} & \\
{\verb|M-x append-to-register|} & \\
{\verb|M-x prepend-to-register|} & \\
{\color{darkred}\verb|M-x dos2unix|} & \\
{\color{darkred}\verb|M-x unix2dos|} & \\
{\color{darkred}\verb|M-x swap-buffers-in-windows|} & \\
{\color{darkred}\verb|M-x xah-show-kill-ring|} & \\
{\verb|M-x customixe-face|} & \\ %% non me lo ricordo mai!
{\verb|M-x load-library|} & load LISP library (from l. path)\\ %% verificare esattamente che significa...
{\verb|M-x eval-expression|} & in minibuffer (alias \verb|M-:|)\\
{\verb|M-x re-builder|} & build regexp interactively\\ %% in Emacs with batteries included 
{\verb|M-x reverse-region|} & reverse lines in region\\
{\verb|M-x untabify|} & replace tabs by spaces\\
{\verb|M-x count-words RET|} & count words in buffer\\
{\verb|M-x man|} & read man pages in emacs\\
{\verb|M-x linum-mode|}  & enable line numbers\\ %% enable line numbers in the margin of the active buffer
{\verb|M-x global-linum-mode|} & enable line numbers globally \\ %% to enable it globally for all buffers run
\end{tabular}

\begin{tabular}{@{}ll@{}}
{\verb|M-x view-mode|} & pager-like for read-only buffers)\\ %% in Emacs with batteries included
{\verb|M-x scroll-lock-mode|} & move buffer instead of the cursor\\ % when navigating %% in Emacs with batteries included
{\verb|M-x follow-mode|} & shows buffer across multiple windows\\ %% in Emacs with batteries included
\verb|M-x compare-windows|  & find next place where text\\ 
& in selected window does not match\\
&text in next window\\
{\verb|M-x occur-edit-mode|} & edit an occur buffer\\ % cfr. Mastering Emacs
{\verb|M-x occur-cease-edit|} & stop editing an ccur buffer\\ % cfr. Mastering Emacs
{\verb|M-x multi-occur|} & search the buffers you specify\\
\end{tabular}

\begin{tabular}{@{}ll@{}}
{\verb|M-x multi-occur-in-matching-buffers|}	& search all buffers that \\
												& match a regexp pattern\\ % cfr. Mastering Emacs
\end{tabular}

%\section{to assigne keybindings functions}
%
%\begin{tabular}{@{}ll@{}}
%{\verb|M-x visual-line-mode|} & activate/deactivate word wrap\\ %% che tasti gli diamo
%\end{tabular}




\vfill\null

%\begin{tabular}{@{}ll@{}}
%
%\end{tabular}

%%%%%%%%%%%%%%%%%%%%%%%%%%%%%%%%%%%%%%%%%%%%%%%%%%%%%%%%%%%%%%
\pagebreak
%\clearpage
\newpage
%%%%%%%%%%%%%%%%%%%%%%%%%%%%%%%%%%%%%%%%%%%%%%%%%%%%%%%%%%%%%%

%%%%%%%%%%%%%%%%%%%%%%%%%%%%%%%
\end{multicols}
\end{document}
%%%%%%%%%%%%%%%%%%%%%%%%%%%%%%%
